\section{The certified positive real period: complete four-label return and correction}
This section applies the full RPCJ calculation to the actual original
real-period point established by the accompanying contraction proof.
The ordered ES--Fable labels and polynomial retain FC5--6, FC15a and
FC31--35 and the original 488-page reader's provenance
\cite{ES488,Cumulative}; the conductor and every norm retain WCF1--6
\cite{WCF}. The conjugation and real-axis positivity used here are
the existing OFM8--11 and earlier PCL15--16 identities
\cite{PCLOriginal,Cumulative}, not a new parity or positivity theorem.

The exact point $(z_*,x_*)$ is the unique real solution of
$f_1(0,z,x)=0$ in the closed rectangle of radius $10^{-15}$ in
each coordinate about the exact rational center
\[
 \begin{split}
 z_c&=0.02211947780866184513200791639903555836071244787217846841,\\
 x_c&=-1.466817440463494183975448871845763045856183562396293859.
 \end{split}\tag{RPE1}
\]
The original parameters are
$\delta=1/4$, $\gamma=3$, $\beta=143/8$, $\eta=21025/256$;
the exact unit is $a=\alpha(1+ix_*)$ for its retained real nonzero
scale $\alpha$, and the original period is $u_*=1/z_*>0$.
The scale cancels with its inverse only in the original conjugated
operator, as in PCL9. Every remaining unit entry is kept there.
The full-series certificate proves $f_1=f_4=0$,
$f_2,f_3\ne0$, and nonzero $z$ and $\xi$ derivatives of $f_1$.
Thus RPCJ8 gives exactly $\nu=2$ and $v=2$ at this point.
The same certificate proves the five-orbit property at this point.
The two vanishing factors are nonzero functions of $\xi$ because
their first $\xi$ derivatives are nonzero; the other two factors
have nonzero value. If a nontrivial conjugated quadric were associated
to $Q_0$, its restriction to the original exponential curve would
vanish identically, contrary to those statements. Association of two
distinct orbit members would reduce to such an association by their
inverse group action. Hence all five original quadrics are distinct.

All moments used below are the actual ones at this certified point.
For a completely explicit finite expression, let $M_j=M_j(z_*,x_*)$
be the original matrix, $v_{j,n}=M_j(\rho_a^n)_a$, and use the
literal polar form
$\mathcal B(v,w)=v_1w_4+v_4w_1-v_2w_3-v_3w_2$.
Then
\[
 a_{j,n}=\frac12\sum_{r=0}^n\binom nr
            \mathcal B(v_{j,r},v_{j,n-r}),\quad
 G_n=\sum_{r=0}^n\binom nr a_{1,r}a_{2,n-r},\quad
 \mu_n=\sum_{r=0}^n\binom nrG_r\bar G_{n-r}.
 \tag{RPE2}
\]
The factor $1/2$ retains the original $Q_0(v)=\mathcal B(v,v)/2$.
These formulas evaluate every needed derivative of the full original
period and factor matrices, including all four multiplicities through
their exact conjugation. The certificate supplies the real enclosures
\[
 \begin{gathered}
 \mu_0=\mu_1=0,\qquad
 \mu_2\in[69.926906\mathbin\pm8.99\,10^{-7}],\\
 \mu_3\in[-11161.6109\mathbin\pm8.96\,10^{-5}],\quad
 \mu_4\in[576883.78\mathbin\pm6.81\,10^{-3}],\quad
 \mu_5\in[5.254785\,10^6\mathbin\pm0.252].
 \end{gathered}\tag{RPE3}
\]
The zero values here are exact consequences of the certified root
and RPCJ3, not the rounding of an interval through zero.
The real parts of the other enclosures are valid because RPCJ4
proves the exact moments real at this real period.

\subsection{Every inverse constant at the certified period}
Retain RPCJ11's exact constants
$c=|G_z|^2$, $R=2\operatorname{Re}(G_z\overline{G_\xi})$,
$d=-iR$, $e=-\mu_2/2$, and
$f=i(\mu_3/6-2\mu_2)$.
The full-series certificate gives
\[
 \begin{gathered}
 c\in[82176.45\mathbin\pm0.00526],\quad
 R\in[3382.6084\mathbin\pm0.0000643],\\
 2c|G_\xi|^2-R^2\in[-5695695\mathbin\pm0.547],\quad
 |\Delta|=|d(d^2-2ce)|\in[1.926631\,10^{10}\mathbin\pm6470].
 \end{gathered}\tag{RPE4}
\]
In particular $\Delta\ne0$ at this actual point; RPCJ22's generic
case therefore occurs. For every retained original $s$, put
\[
 C_s=\begin{pmatrix}
 -\frac{\mu_2}2\frac{\rho_2}{\rho_0}&
 i(\mu_3/6-2\mu_2)\frac{\rho_3}{\rho_0}\\
 0&-\frac{\mu_2}2\frac{\rho_3}{\rho_1}
 \end{pmatrix},\qquad
 (p_1,p_2,p_3,p_4)=(5,8,8,8).
 \tag{RPE5}
\]
The complete exact exterior constants are
\[
 \lim_{z\to z_*}|z-z_*|^{p_r}\|\bigwedge^rB_{3,s}(z)^{-1}\|
 =\left(
 \frac{|\Delta|\rho_3}{c^4\rho_0},\quad
 \frac{\mu_2^2\rho_2\rho_3}{4c^4\rho_0\rho_1},\quad
 \frac{\|C_s\|}{c^4},\quad c^{-4}
 \right)_r.
 \tag{RPE6}
\]
The block norm is the explicit expression RPCJ21 applied to the
three entries in RPE5. This retains the actual $\mu_3$ rather than
bounding or deleting its off-diagonal contribution.
The inverse singular exponents are $(5,3,0,0)$ and their constants
are exactly
\[
 \left(
 \frac{|\Delta|\rho_3}{c^4\rho_0},\quad
 \frac{\mu_2^2\rho_2}{4|\Delta|\rho_1},\quad
 \sigma_-(C_s)^{-1},\quad\sigma_+(C_s)^{-1}
 \right).
 \tag{RPE7}
\]
At $s=1$, the original weights give
$(\rho_0,\rho_1,\rho_2,\rho_3)=(1,\sqrt2,\sqrt{8/3},\sqrt{16/5})$.
Propagation of the saved certificate enclosures through these exact
formulas gives all four exterior constants, respectively,
\[
 \begin{split}
 &[7.55760\,10^{-10}\mathbin\pm3.35\,10^{-16}],\quad
 [5.53711\,10^{-17}\mathbin\pm5.87\,10^{-23}],\\
 &[7.84748\,10^{-17}\mathbin\pm6.86\,10^{-23}],\quad
 [2.192858\,10^{-20}\mathbin\pm6.28\,10^{-27}].
 \end{split}\tag{RPE8}
\]
The last two inverse singular limits at this order are in
\[
[1.4172517\mathbin\pm7.40\,10^{-8}],\qquad
[0.00027943449\mathbin\pm5.32\,10^{-12}].
\]
The second inverse
singular leading constant is in
$[7.32655\,10^{-8}\mathbin\pm1.65\,10^{-14}]$.
These bounds are propagated from the full original certificate,
not new finite-series substitutions for the actual functions.

The actual original period satisfies
$u_*\in[45.20902385898\mathbin\pm3.31\,10^{-12}]$.
In the $u$ coordinate, multiply the four constants in RPE6 by
$u_*^{10},u_*^{16},u_*^{16},u_*^{16}$ respectively, from
$z-z_*=-(u-u_*)/(uu_*)$.
At $s=1$ they are enclosed respectively by
\[
 [26954700\mathbin\pm55.4],\quad
 [1.686109\,10^{10}\mathbin\pm7850],\quad
 [2.38964\,10^{10}\mathbin\pm10900],\quad
 [6677480\mathbin\pm2.02].
 \tag{RPE9}
\]

\subsection{All four original labels on the old and corrected source spaces}
The literal original marking remains
\[
 \begin{array}{c|cccc}
 a&\mathrm M&\mathrm N&\mathrm T&\mathrm E\\\hline
 \rho_a&3/4+3i&3/4-3i&1/4+3i&1/4-3i
 \end{array},\qquad
 \ell_a(S)=\prod_{b\ne a}\frac{S-\rho_b}{\rho_a-\rho_b}
             =\sum_{r=0}^3\lambda_{ra}S^r.
 \tag{RPE10}
\]
Define $L=(\lambda_{ra})=(V_\rho^{\mathsf T})^{-1}$, as in FC26--35.
Every coefficient is specified by the displayed product; in particular
\[
 \begin{gathered}
 (\lambda_{3\mathrm M},\lambda_{3\mathrm N},
   \lambda_{3\mathrm T},\lambda_{3\mathrm E})
 =\frac1{435}(-24-2i,-24+2i,24-2i,24+2i),\\
 \lambda_{2a}=(\rho_a-2)\lambda_{3a},\quad
 \lambda_{1a}=(\rho_a^2-2\rho_a+155/8)\lambda_{3a},\\
 \lambda_{0a}=(\rho_a^3-2\rho_a^2+(155/8)\rho_a-147/8)\lambda_{3a}.
 \end{gathered}\tag{RPE11}
\]
These formulas follow by dividing the full root polynomial
$\prod_a(S-\rho_a)$ by its factor $S-\rho_a$ and retaining the
original derivative denominator.

On the old degree-three source, the four outputs at this period are
\[
 y_a^{\rm old}=T_A\ell_a
       =\lambda_{2a}\mu_2+\lambda_{3a}(\mu_3+3\mu_2S'),\qquad
 Y^{\rm old}=
 \begin{pmatrix}
 0&0&\mu_2&\mu_3\\0&0&0&3\mu_2\\0&0&0&0\\0&0&0&0
 \end{pmatrix}L.
 \tag{RPE12}
\]
The rank is exactly two and the kernel on that original source is
$\mathcal P_1$. In label coordinates this kernel is exactly
$\operatorname{span}\{\boldsymbol1,\rho\}$, since interpolation
sends those two evaluation vectors to $1,S$.
In the original ES coordinates the kernel of
$Y^{\rm old}K^{-1}$ is therefore
$K\operatorname{span}\{\boldsymbol1,\rho\}$, with the full
unchanged $K$ of FC15a. No four-label column is removed or identified.

The actual order-two source, with its complete original quotient norm,
is instead $\mathcal U_2=\mathcal P_5/\mathcal P_1$.
Its four retained inputs and outputs are
\[
 x_a^{(2)}=[S^2\ell_a(S)],\qquad
 y_a^{(2)}=T_A(S^2\ell_a),\qquad
 Y^{(2)}=
 \begin{pmatrix}
 \mu_2&\mu_3&\mu_4&\mu_5\\
 0&3\mu_2&4\mu_3&5\mu_4\\
 0&0&6\mu_2&10\mu_3\\
 0&0&0&10\mu_2
 \end{pmatrix}L.
 \tag{RPE13}
\]
For its proof use the complete original Taylor translation formula
\[T_AS^{2+r}=\sum_{k=0}^r\binom{2+r}{k}\mu_{2+r-k}S'^k.\]
This evaluates every entry with the actual moments RPE2--3.
Here
\[
 \det V_\rho^{\mathsf T}=-1305/4,\qquad
 \det Y^{(2)}=\frac{180\mu_2^4}{-1305/4}
       =-\frac{16}{29}\mu_2^4
       \in[-13191650\mathbin\pm4.20]\ne0.
 \tag{RPE14}
\]
Thus the corrected four-column map is exactly invertible, with all
four original labels and the actual positive $\mu_2$.

There is also an exact operator identity giving every correction to
the old map in RPE12. Define the unchanged derived symbols
$T_A^{[r]}P(S')=\sum_{u,w}a_{uw}\beta_{8,u,w}^{r}
P(S'+\beta_{8,u,w})$, for $r=1,2$.
Expansion of $(S'+\beta_{8,u,w})^2$ inside the original sum proves
\[
 T_A(S^2P)=S'^2T_AP+2S'T_A^{[1]}P+T_A^{[2]}P.
 \tag{RPE15}
\]
It retains all original coefficient and frequency contributions.
The complete matrices of the last two maps on the old monomial
source basis $(1,S,S^2,S^3)$ are
\[
 T_A^{[1]}=
 \begin{pmatrix}0&\mu_2&\mu_3&\mu_4\\
 0&0&2\mu_2&3\mu_3\\0&0&0&3\mu_2\\0&0&0&0\end{pmatrix},
 \quad
 T_A^{[2]}=
 \begin{pmatrix}\mu_2&\mu_3&\mu_4&\mu_5\\
 0&\mu_2&2\mu_3&3\mu_4\\0&0&\mu_2&3\mu_3\\0&0&0&\mu_2\end{pmatrix}.
 \tag{RPE16}
\]
Multiplying their coefficient matrices by $L$ and applying RPE15
recovers RPE13 entry by entry. In particular the lost two old
directions are restored by specified original moment terms, not by
discarding the original kernel or choosing a different amplitude.

The source map $P\mapsto[S^2P]$ is an isomorphism of the underlying
four-dimensional spaces, with inverse obtained from the unique
representative containing only degrees $2$ through $5$ and division
by $S^2$. Its exact norm is the pullback of the stated Gamma quotient
norm; it is not asserted isometric to the old degree-three norm.
Writing $h_{n,s}=M_sn!(s/2)_n$, the full Gram determinants are
\[
 \det G_{x,2}=
   \frac{h_{2,s}h_{3,s}h_{4,s}h_{5,s}}{(1305/4)^2},\qquad
 \det G_{y,2}=
   \frac{180^2\mu_2^8h_{0,s}h_{1,s}h_{2,s}h_{3,s}}{(1305/4)^2},
 \quad
 \frac{\det G_{y,2}}{\det G_{x,2}}
   =\left(\frac{\mu_2}2\right)^8
      \prod_{r=0}^3\frac{\rho_{r+2}^2}{\rho_r^2}.
 \tag{RPE17}
\]
These follow from the unchanged leading coefficients of the original
orthonormal polynomials, exactly as in FC28--30. Both Gram
determinants retain the entire mass $M_s^4$.

Finally define $X_2e_a=x_a^{(2)}$, $Y_2e_a=y_a^{(2)}$ and retain
$\Phi_2=X_2K^{-1}$, $\Psi_2=Y_2K^{-1}$.
The full original four-label intertwining is
\[
 T_A\Phi_2=\Psi_2,\qquad
 P_{\mathcal U_2}=\Phi_2P\Phi_2^{-1},\quad
 P_{\mathcal W}=\Psi_2P\Psi_2^{-1},\qquad
 T_AP_{\mathcal U_2}=P_{\mathcal W}T_A,
 \tag{RPE18}
\]
with the complete original $P$ from FC31, and
$P_{\mathcal W}(y_a^{(2)})=2i\,y_{\mathrm M}^{(2)}$ for every
$a=\mathrm M,\mathrm N,\mathrm T,\mathrm E$.
Its inverse matrices are the adjugates divided by RPE14 and FC15a's
nonzero determinants; the pulled-back metrics are
$(K^{-1})^*G_{x,2}K^{-1}$ and $(K^{-1})^*G_{y,2}K^{-1}$.
Thus the actual real-period degeneration, its full fixed-space poles,
and its exact order-two four-label correction are retained together.

\subsection{The two lost scalars, their complete metric, and the total return}
The rank-two old receiver has four distinct marked outputs. Indeed,
in the ordered pair list $(\mathrm{MN},\mathrm{MT},\mathrm{ME},
\mathrm{NT},\mathrm{NE},\mathrm{TE})$, RPE11 gives
\[
 \lambda_{3a}-\lambda_{3b}
  =\frac1{435}(-4i,-48,-48-4i,-48+4i,-48,-4i)_{ab}.
 \tag{RPE19}
\]
The coefficient of $S'$ in $y_a^{\rm old}-y_b^{\rm old}$ is
$3\mu_2(\lambda_{3a}-\lambda_{3b})\ne0$. This proves all six
distinctions without identifying a label with a conjugate label.
For the exact metric statement, put $b=s/2$, $M=(2\pi)^{s/2}$,
$c=a_{\rm amp}/2$ and $c'=c-4$. If this difference is $d_0+d_1S'$,
its original squared norm is
$M(|d_0+c'd_1|^2+b|d_1|^2)>0$.

Let $\Phi_0q=\sum_a(K^{-1}q)_a\ell_a(S)$ and
$J_0=T_A\Phi_0$. Retain the literal two coefficient functionals
$t_r(q)=[S^r]\Phi_0q$, $r=0,1$. Define the complete linear map
\[
 \Theta_2:\mathbb C^4\longrightarrow\mathcal P_1\oplus\mathbb C^2,
 \qquad \Theta_2q=(J_0q,t_0(q),t_1(q)).
 \tag{RPE20}
\]
Here the first $\mathcal P_1$ uses the original target variable $S'$;
the two scalar coordinates are the coefficients lost by the original
degree-three conductor. For $Y=Y_0+Y_1S'$, put
\[
 \begin{gathered}
 c_3=\frac{Y_1}{3\mu_2},\qquad
 c_2=\frac{Y_0-\mu_3c_3}{\mu_2},\qquad
 p_Y(S)=c_2S^2+c_3S^3,\\
 \Theta_2^{-1}(Y,t_0,t_1)
       =K\bigl(p_Y(\rho_a)+t_0+t_1\rho_a\bigr)_a.
 \end{gathered}\tag{RPE21}
\]
The two coefficients of $T_Ap_Y$ are precisely $Y_0,Y_1$ by
RPE12, whereas $p_Y$ itself has zero constant and linear coefficients.
These identities prove $\Theta_2\Theta_2^{-1}=I$.
Conversely any polynomial in the old source differs from its
corresponding $p_Y$ by exactly its original constant and linear
terms; interpolation proves the other composition. In coefficient
bases the determinant of $\Theta_2$ is
$3\mu_2^2/((\det V_\rho^{\mathsf T})(\det K))\ne0$.
Thus both lost scalars, including their original ES coordinates,
are recovered by finite displayed formulas.

Every cross term of the original source metric is as follows. Set
$m_4=b(3b+2)$ and $m_6=b(15b^2+30b+16)$, and define
\[
 \begin{gathered}
 Q=\begin{pmatrix}1&c\\c&c^2+b\end{pmatrix},\qquad
 U=\begin{pmatrix}
 c^2-b&c^3-3cb\\c^3+cb&c^4-m_4
 \end{pmatrix},\\
 H=\begin{pmatrix}
 c^4+2c^2b+m_4&c(c^4+2c^2b+m_4)\\
 c(c^4+2c^2b+m_4)&c^6+3c^4b+3c^2m_4+m_6
 \end{pmatrix}.
 \end{gathered}\tag{RPE22}
\]
With $v=(c_2,c_3)^{\mathsf T}$ and $t=(t_0,t_1)^{\mathsf T}$,
direct integration against the unchanged WCF2 measure gives
\[
 \begin{split}
 \|p_Y+t_0+t_1S\|_{s,c}^2
 &=M\{v^*Hv+2\operatorname{Re}(t^*Uv)+t^*Qt\}\\
 &=M v^*(H-U^*Q^{-1}U)v
       +M(t+Q^{-1}Uv)^*Q(t+Q^{-1}Uv),\\
 Q^{-1}Uv
 &=\begin{pmatrix}
 -(c^2+b)c_2+(-2c^3+2c)c_3\\
 2cc_2+(3c^2-3b-2)c_3
 \end{pmatrix},\\
 H-U^*Q^{-1}U
 &=\begin{pmatrix}
 2b(b+1)&6cb(b+1)\\
 6cb(b+1)&18c^2b(b+1)+6b(b+1)(b+2)
 \end{pmatrix}.
 \end{split}\tag{RPE23}
\]
For completeness, WCF2's recurrence gives the unchanged measure's
moments $M$, $Mb$, $Mm_4$, $Mm_6$ in even degrees zero through six.
One may verify them from
$p_2=y^2-b$, $p_3=y^3-(3b+2)y$ and their exact WCF2 norms.
Odd moments vanish. Substitution of $S=c+iy$ then gives every
entry of RPE22; for example
$M^{-1}\langle S,S^3\rangle=c^4-m_4$ and
$M^{-1}\langle S^2,S^3\rangle=c(c^4+2c^2b+m_4)$.
The inverse of $Q$ is
$b^{-1}\left(\begin{smallmatrix}c^2+b&-c\\-c&1\end{smallmatrix}\right)$.
Multiplication gives the last two expressions in RPE23 and proves
the completed square. The original quotient norm is its first term,
attained at the exact two-scalar value $t=-Q^{-1}Uv$.
No source mass, center, kernel direction or cross term is omitted.

For a full Gram directly in the displayed augmented coordinates,
define
\[
 D_Y=\begin{pmatrix}
 1/\mu_2&-\mu_3/(3\mu_2^2)\\0&1/(3\mu_2)
 \end{pmatrix},\qquad
 G_{\Theta_2}=M\begin{pmatrix}
 D_Y^*HD_Y&D_Y^*U^*\\UD_Y&Q
 \end{pmatrix}.
 \tag{RPE24}
\]
Thus $v=D_Y(Y_0,Y_1)^{\mathsf T}$ and RPE24 is exactly the
original source metric transported by RPE21, in the order
$(Y_0,Y_1,t_0,t_1)$. Its quotient block is
$M D_Y^*(H-U^*Q^{-1}U)D_Y$.
The target norm stays
$\|Y\|_{s,c'}^2=M(|Y_0+c'Y_1|^2+b|Y_1|^2)$; it is not
replaced by that source quotient metric.

The four marked total states and their full common images retain
both scalars:
\[
 \begin{gathered}
 \Theta_2(Ke_a)=(y_a^{\rm old},\lambda_{0a},\lambda_{1a}),\\
 \Theta_2P(Ke_a)
   =(2i y_{\mathrm M}^{\rm old},
          2i\lambda_{0\mathrm M},2i\lambda_{1\mathrm M}).
 \end{gathered}\tag{RPE25}
\]
All constants on the right are RPE11--12's exact original values.
Finally the entire polynomial return is
\[
 \widehat P_2(Y,t_0,t_1)
  =\Theta_2P\!\left(K(p_Y(\rho_a)+t_0+t_1\rho_a)_a\right),
 \quad \Theta_2P=\widehat P_2\Theta_2,
 \quad\det D\widehat P_2=-2.
 \tag{RPE26}
\]
The full original $P$ is FC31, so these are explicit finite
polynomial expressions with inverse coordinate map RPE21.
The chain rule proves the determinant because the two constant
determinants of $\Theta_2$ and its inverse cancel. This uses the
existing FC31 Jacobian identity with its ES--Fable provenance; it
asserts no new counterexample beyond that source polynomial.
Equations RPE20--26 retain the old map with both lost scalars;
RPE13--18 separately give the actual order-two conductor on its
correct quotient. The exact morphisms and their different original
metrics have both been computed.
No assertion places the chosen geometric quartet or unit on the
arithmetic Riemann-zeta locus.
